% !TeX program = xelatex
% Bao cao nhom 10 - Inside Airbnb Toronto - ban du thao gan hoan thien.
% Bien dich: XeLaTeX (Overleaf: Menu -> Compiler -> XeLaTeX).
% Cac bieu do duoc ve truc tiep bang PGFPlots, khong can tep hinh ngoai.
% Truoc khi nop: hoan tat cac muc \pending{...}, xac nhan ngay mốc dữ liệu
% 2025-12-12 voi giang vien, chay Gemini bộ nhớ đệm/evaluation va cap nhat commit.

\documentclass[11pt,a4paper]{article}

% --- Ngon ngu va font ---
\usepackage{fontspec}
\IfFontExistsTF{Times New Roman}{
  \setmainfont{Times New Roman}
}{
  \IfFontExistsTF{TeX Gyre Termes}{
    \setmainfont{TeX Gyre Termes}
  }{
    \setmainfont{Noto Serif}
  }
}
\usepackage[vietnamese]{babel}

% --- Bo cuc va goi ho tro ---
\usepackage[a4paper,margin=2.15cm,headheight=14pt]{geometry}
\usepackage{amsmath,amssymb}
\usepackage{graphicx}
\usepackage{booktabs,tabularx,array,longtable}
\usepackage{caption}
\usepackage{enumitem}
\usepackage{xcolor}
\usepackage{hyperref}
\usepackage{needspace}
\usepackage{pgfplots}
\usetikzlibrary{positioning,arrows.meta}
\pgfplotsset{compat=1.18}

\hypersetup{
  colorlinks=true,
  linkcolor=blue!45!black,
  urlcolor=blue!55!black,
  citecolor=blue!45!black,
  pdftitle={Phân tích nhóm mẫu Airbnb tại Toronto - Nhóm 10},
  pdfauthor={Nhom 10 - UET.DSE2049}
}
\setlength{\parindent}{1em}
\setlength{\parskip}{2pt}
\setlength{\emergencystretch}{2em}
\setlist{nosep,leftmargin=*}
\captionsetup{font=small,labelfont=bf,skip=3pt}
\renewcommand{\figurename}{Hình}
\renewcommand{\tablename}{Bảng}
\renewcommand{\abstractname}{Tóm tắt}
\renewcommand{\refname}{Tài liệu tham khảo}
\newcolumntype{Y}{>{\raggedright\arraybackslash}X}
\newcolumntype{C}{>{\centering\arraybackslash}X}
\newcommand{\code}[1]{\texttt{\small #1}}
\newcommand{\pct}[1]{#1\%}
\newif\ifdraft
\drafttrue
\newcommand{\pending}[1]{\ifdraft\textcolor{red!70!black}{\textbf{[CẦN HOÀN THIỆN: #1]}}\else\textbf{[CẦN ĐIỀN]}\fi}
\newcommand{\snapshot}[1]{\texttt{#1}}

\title{\vspace{-1.5em}\Large\bfseries Phân tích nhóm mẫu chỗ ở Airbnb tại Toronto:\\[-2pt]
\large Nguồn cung, chất lượng dữ liệu, lịch sẵn có và phản hồi khách lưu trú}
\author{\textbf{Nhóm 10}\\[-1pt]
\small Dương Nhật Huy \quad Mai Anh Đức \quad Nguyễn Quốc Hiếu\\[-1pt]
\small Vũ Phạm Khánh Nam \quad Phan Ngọc Huy\\[-1pt]
\small Lập trình xử lý dữ liệu (UET.DSE2049 1) -- Học kỳ I, 2026--2027\\[-1pt]
\small Viện Trí tuệ nhân tạo, Trường Đại học Công nghệ, ĐHQGHN\\[-1pt]
\small MSSV: \pending{bổ sung MSSV của 5 thành viên}}
\date{}

\begin{document}
\maketitle
\vspace{-1.6em}
\noindent\small
\textbf{Bên nhận phân tích giả định:} Nhóm theo dõi thị trường lưu trú ngắn hạn và du lịch Toronto.\\
\textbf{Mã dữ liệu:} \code{DSE2049-C01-G10-toronto} \quad
\textbf{Tiền tệ:} CAD \quad
\textbf{Thành phố:} Toronto, Canada.\\
\textbf{Kho mã:} \href{https://github.com/uet-iai-notebook-labs-2026/2627-LTXLDL-Airbnb-dse2049_1-nhom10}{GitHub -- Nhóm 10}.\\
\textbf{Commit mã nguồn dùng để đối chiếu:} \code{d045d55869b76c21b92b8b48bf2f0417a929efca} (nhánh \code{main}; chưa phải tag \code{final}).\\
\textbf{Trạng thái:} Bản thảo gần hoàn thiện; kết quả Gemini và xác nhận về một mốc dữ liệu vẫn còn chờ.
\par\normalsize

\begin{abstract}
Báo cáo phân tích 1.000 mã chỗ ở Airbnb được giao cho Nhóm 10 tại Toronto qua bốn mốc dữ liệu từ tháng 9/2025 đến tháng 6/2026. Nhóm kiểm tra 12 tệp dữ liệu dạng bảng, xây dựng quy trình xử lý có thể tái lập và kiểm thử tích hợp trên dữ liệu thật ở chế độ không sử dụng mô hình ngôn ngữ lớn. Cả bốn mốc đều ghi nhận đủ 1.000 chỗ ở; tại mốc cuối có 980 chỗ ở có giá hợp lệ, với giá trung vị 213,55 CAD/đêm. Trong nhóm mẫu, 87,5\% chỗ ở là toàn bộ nhà/căn hộ; 11,6\% chỗ ở thuộc về chủ nhà có nhiều hơn một chỗ ở được quan sát. Mã chủ nhà (\code{host\_id}) thay đổi giữa các mốc liền kề lần lượt là 3, 35 và 37 trên 1.000 cặp. Độ phủ dữ liệu lịch trong cửa sổ 90 ngày đạt từ 99,38\% đến 100\%. Bộ 100 nhãn cảm xúc được gán thủ công đã hoàn tất, nhưng bộ nhớ đệm Gemini cho 1.000 bài đánh giá chưa có; vì vậy báo cáo chưa công bố độ chính xác của mô hình hoặc kết luận về các khía cạnh trải nghiệm khách. Các kết quả chỉ mô tả nhóm mẫu được giao, không đại diện cho toàn bộ thị trường Toronto.
\end{abstract}
\noindent\textbf{Từ khóa:} Inside Airbnb; Toronto; nhóm mẫu; kiểm soát chất lượng dữ liệu; lịch sẵn có; phân tích đánh giá; Gemini.

\section{Giới thiệu và bài toán}
Airbnb cung cấp dữ liệu quan sát được về các chỗ ở được niêm yết, lịch sẵn có và phản hồi sau lưu trú. Tuy nhiên, dữ liệu nhiều bảng và nhiều thời điểm có thể khác nhau về cấu trúc cột, mức độ thiếu và khóa liên kết. Nếu không kiểm tra trước khi tổng hợp, số lượng chỗ ở, giá trung vị, tỷ lệ sẵn có hoặc chỉ số theo vùng có thể bị sai lệch.

Báo cáo hướng tới một bên nhận phân tích giả định là nhóm theo dõi thị trường lưu trú ngắn hạn và du lịch Toronto. Mục tiêu là chuyển dữ liệu được giao thành các chỉ số có định nghĩa rõ, có thể truy vết về mã và dữ liệu đầu vào, qua đó hỗ trợ đặt câu hỏi tiếp theo về cơ cấu nguồn cung, lịch sẵn có và trải nghiệm khách.

Phạm vi nghiên cứu là nhóm mẫu cố định gồm 1.000 \code{listing\_id} do bảng tính phân công của học phần quy định, không phải toàn bộ chỗ ở tại Toronto. Do nhóm mẫu được cố định theo danh sách mã, việc một mã xuất hiện hoặc vắng mặt trong bảng đầu vào chỉ phản ánh độ phủ dữ liệu; không được diễn giải thành chỗ ở mới gia nhập hoặc rời khỏi thị trường. Tương tự, ngày lịch không còn sẵn không đồng nghĩa chắc chắn đã có khách đặt.

\section{Câu hỏi phân tích và thiết kế nghiên cứu}
Nghiên cứu tổ chức theo ba câu hỏi:

\noindent\textbf{RQ1 -- Không gian và phân khúc.} Nhóm mẫu được phân bố như thế nào giữa tám vùng phân tích của bảng tính, và cơ cấu loại hình chỗ ở ra sao? Giá niêm yết ở mốc dữ liệu cuối cho thấy mức giá điển hình nào trong phạm vi mẫu?

\noindent\textbf{RQ2 -- Biến động theo thời gian.} Độ phủ nhóm mẫu, các thuộc tính chỗ ở và mức độ đầy đủ của dữ liệu lịch thay đổi như thế nào qua bốn mốc dữ liệu? Những thay đổi quan sát được có đủ căn cứ để kết luận về biến động của toàn thị trường hay không?

\noindent\textbf{RQ3 -- Phản hồi khách lưu trú.} Các bài đánh giá được giao có phân bố nhãn cảm xúc thủ công như thế nào, và mô hình Gemini có phân loại đáng tin cậy hơn phương pháp đường cơ sở dùng từ khóa hay không? Chỉ kết luận câu hỏi thứ hai sau khi tạo bộ nhớ đệm thật và đánh giá hai phương pháp trên cùng tập 100 nhãn tay.

\section{Dữ liệu và phương pháp xử lý}
\subsection{Nguồn dữ liệu và phạm vi}
Nguồn dữ liệu là Inside Airbnb, được phát hành theo giấy phép CC BY 4.0 \cite{inside,snap1,snap2,snap3,snap4}. Theo tệp kê khai phân công của nhóm, bốn mốc dữ liệu được cấu hình là 2025-09-11, 2025-12-12, 2026-03-19 và 2026-06-15. Với mỗi mốc, nhóm kiểm tra một tệp thông tin chỗ ở (listings), một tệp lịch (calendar) và một tệp bài đánh giá (reviews); cả 12/12 tệp dạng bảng này được báo cáo đọc thành công. Tệp GeoJSON được lưu riêng để hỗ trợ đối chiếu không gian; tài liệu kiểm chứng hiện chỉ xác nhận việc đối chiếu tệp này ở mốc 2025-09-11. Quy mô dữ liệu theo từng mốc được tóm tắt tại Bảng~\ref{tab:data}.

\begin{table}[htbp]
\centering\small
\caption{Quy mô dữ liệu đầu vào theo các mốc được giao.}\label{tab:data}
\begin{tabularx}{\linewidth}{@{}lCCCCY@{}}
\toprule
Mốc dữ liệu & Chỗ ở & Lịch & Bài đánh giá & GeoJSON & Ghi chú\\
\midrule
\snapshot{2025-09-11} & 1.000 & 365.000 & 5.000 & 1 tệp & Đọc thành công\\
\snapshot{2025-12-12} & 1.000 & 365.000 & 5.000 & 1 tệp & Ngày cần xác nhận\\
\snapshot{2026-03-19} & 1.000 & 365.000 & 5.000 & 1 tệp & Đọc thành công\\
\snapshot{2026-06-15} & 1.000 & 365.000 & 5.000 & 1 tệp & Mốc tính giá\\
\bottomrule
\end{tabularx}
\end{table}

Bảng thông tin chỗ ở (listings) và bảng lịch (calendar) được lọc theo 1.000 \code{listing\_id} trong bảng tính phân công. Bảng bài đánh giá (reviews) được lọc độc lập theo 5.000 \code{review\_id} được giao ở mỗi mốc, không lọc thêm theo \code{listing\_id}; sau khi loại các bài đánh giá trùng giữa các mốc, quy trình xử lý ghi nhận 5.000 bài đánh giá duy nhất. Trong đó, 1.000 bài thuộc mẫu phân tích bằng mô hình ngôn ngữ lớn và 100 bài thuộc tập gán nhãn thủ công. Mã định danh được giữ dưới dạng chuỗi để tránh mất độ chính xác với các mã dài.

\textbf{Lưu ý về ngày mốc dữ liệu.} Tệp kê khai và bảng tính phân công của nhóm ghi ngày thứ hai là \snapshot{2025-12-12}, trong khi trang đề bài hiện được nhóm ghi nhận là \snapshot{2025-11-11}. Nhóm giữ nguyên ngày trong bảng tính phân công và dữ liệu đã tải để bảo toàn phạm vi được giao; \pending{nhóm trưởng xác nhận với giảng viên ngày đúng và cập nhật nếu cần}. Không tự thay ngày mốc trong báo cáo nếu chưa có xác nhận.

\subsection{Cấu trúc và liên kết bảng}
Một dòng trong bảng thông tin chỗ ở (listings) đại diện cho một chỗ ở tại một mốc dữ liệu, với khóa \code{(snapshot\_date, listing\_id)}. Bảng lịch có đơn vị quan sát là từng chỗ ở--ngày, dùng khóa \code{(snapshot\_date, listing\_id, date)} trước khi tổng hợp. Bảng bài đánh giá (reviews) có khóa \code{review\_id}; bài đánh giá trùng giữa các mốc được loại trùng khi phân tích, đồng thời lưu \code{first\_seen\_snapshot}. GeoJSON cung cấp ranh giới khu vực hành chính để liên kết dữ liệu không gian.

Bảng tính phân công định nghĩa tám vùng chữ nhật \code{analysis\_area} từ A01 đến A08. Đây là các vùng phân tích của bài tập, không phải ranh giới hành chính \code{neighbourhood\_cleansed}. Vì vậy, kết quả theo A01--A08 và kết quả theo khu vực hành chính được giữ thành hai chiều không gian riêng. Quy trình hiện xuất 500 dòng KPI theo khu vực hành chính, tương ứng 125 khu vực qua bốn mốc dữ liệu.

\subsection{Quy trình xử lý có thể tái lập}
Quy trình xử lý chung gồm các bước: đọc cấu hình và tệp kê khai; kiểm tra/tải dữ liệu đầu vào; lọc đúng phạm vi; kiểm soát chất lượng bảng listings; tổng hợp bảng lịch và bảng bài đánh giá; tính KPI; chạy phương pháp đường cơ sở; đọc bộ nhớ đệm Gemini (ở chế độ đầy đủ); đánh giá bằng nhãn thủ công; cuối cùng xuất CSV, siêu dữ liệu và hình. Các mô-đun được tách theo trách nhiệm thành viên trong \code{src/airbnb/}. Dữ liệu gốc và bảng tính phân công không bị sửa. Siêu dữ liệu của mỗi lần chạy lưu cấu hình, mã băm kiểm tra đầu vào, phiên bản Python, số dòng và danh sách tệp kết quả.

\begin{figure}[htbp]
\centering
\begin{tikzpicture}[node distance=4mm and 4mm, every node/.style={font=\small,align=center},
  procbox/.style={draw,rounded corners,fill=blue!4,minimum height=8mm,text width=2.6cm},
  finalbox/.style={draw,rounded corners,fill=green!5,minimum height=8mm,text width=2.7cm},
  flowarrow/.style={->,>=stealth,thick}]
\node[procbox] (src) {Nguồn Inside Airbnb\\+ danh sách mã được giao};
\node[procbox,right=of src] (load) {Đọc, kiểm tra\\và lọc phạm vi};
\node[procbox,right=of load] (qa) {Kiểm soát chất lượng\\bảng chỗ ở và lịch};
\node[procbox,right=of qa] (kpi) {Chỉ số, nhóm mẫu\\và phương pháp cơ sở};
\node[finalbox,right=of kpi] (out) {CSV, siêu dữ liệu\\và hình};
\draw[flowarrow] (src)--(load); \draw[flowarrow] (load)--(qa); \draw[flowarrow] (qa)--(kpi); \draw[flowarrow] (kpi)--(out);
\node[procbox,below=7mm of kpi,text width=4.3cm] (llm) {Bộ nhớ đệm Gemini đã kiểm tra nguồn gốc\\+ đánh giá cùng 100 nhãn tay};
\draw[flowarrow] (qa.south) |- (llm.west);
\draw[flowarrow] (llm.east) -| (out.south);
\end{tikzpicture}
\caption{Kiến trúc xử lý ở mức tổng quan. Nhánh Gemini chỉ được kích hoạt khi bộ nhớ đệm thật có đủ 1.000 bài đánh giá; chế độ không dùng mô hình xuất kết quả vào thư mục riêng.}
\label{fig:pipeline}
\end{figure}

\subsection{Kiểm soát chất lượng dữ liệu}
Nhóm áp dụng nguyên tắc ưu tiên gắn cờ thay vì xóa dữ liệu. Dòng chỉ bị loại khi không có mã chỗ ở hợp lệ hoặc trùng khóa trong cùng một mốc dữ liệu. Giá không hợp lệ được đặt thiếu để không tham gia phép tính giá, nhưng chỗ ở vẫn được giữ trong mẫu số nguồn cung. Cách làm này tránh làm giảm nguồn cung chỉ vì một dòng thiếu giá.

Bộ kiểm soát chất lượng bảng listings có 19 quy tắc, bao gồm thiếu trường, sai kiểu, trùng khóa, ngoại lai, sai ngày, lệch cấu trúc cột, sai danh mục và kiểm tra giữa các mốc dữ liệu. Ngưỡng cấu hình gồm giá tối đa 10.000 CAD/đêm và \code{minimum\_nights} tối đa 365. Đối với nhóm mẫu được giao, bước kiểm soát chất lượng không loại dòng: 4.000/4.000 bản ghi chỗ ở vẫn được giữ; mốc cuối có 20 chỗ ở thiếu giá và 980 chỗ ở có giá hợp lệ. Báo cáo kiểm soát chất lượng ghi nhận 20 dòng bị gắn cờ do thiếu giá ở mốc cuối; không dùng giá thiếu để tính trung vị.

\begin{table}[htbp]
\centering\small
\caption{Các phép kiểm tra trọng yếu và hành động xử lý.}\label{tab:qa}
\begin{tabularx}{\linewidth}{@{}lYY@{}}
\toprule
Nhóm kiểm tra & Điều kiện/ý nghĩa & Hành động và kết quả quan sát\\
\midrule
Khóa và trùng & Mã định danh rỗng, sai định dạng hoặc trùng trong cùng một mốc dữ liệu & Kiểm tra đầu vào không ghi nhận ID thiếu/trùng; không làm giảm quy mô 1.000 chỗ ở/mốc.\\
Trường quan trọng & Thiếu \code{host\_id}, \code{room\_type} hoặc vùng phân tích & Gắn cờ theo trường; không tự điền giá trị không có căn cứ.\\
Giá & Thiếu, không chuyển đổi được sang số, không dương hoặc vượt ngưỡng & Chỉ phân tích giá tại 2026-06-15; 20/1.000 giá thiếu, giữ chỗ ở và loại giá khỏi thống kê.\\
Danh mục và cấu trúc cột & Loại phòng ngoài danh mục; số/tên cột thay đổi & Ghi nhận cấu trúc cột tăng 79 $\rightarrow$ 85 $\rightarrow$ 90 $\rightarrow$ 90 cột; không tự đồng nhất cột khác nghĩa.\\
Ngày và tọa độ & Ngày không hợp lệ, quét lệch, tọa độ ngoài phạm vi hoặc dịch chuyển lớn & Gắn cờ/ghi nhận; không tự sửa tọa độ. Các mốc dữ liệu có độ trễ quét quan sát 0--2 ngày.\\
So sánh liên mốc & Cùng chỗ ở đổi chủ nhà, loại phòng, khu vực hành chính hoặc tọa độ & Báo cáo số cặp thay đổi; không diễn giải tự động thành đổi chủ hay thay đổi thị trường.\\
\bottomrule
\end{tabularx}
\end{table}

Số chỗ ở thiếu giá trong dữ liệu nguồn theo bốn mốc lần lượt là 69/1.000, 1.000/1.000, 28/1.000 và 20/1.000. Do đó nhóm chỉ dùng giá tại mốc cuối theo cấu hình \code{price\_snapshot}. Bước kiểm soát chất lượng không làm thay đổi số lượng chỗ ở hoặc giá trung vị trước và sau xử lý trong nhóm mẫu được báo cáo; giá trung vị ở mốc cuối là khoảng 213,55 CAD/đêm. Đây là giá niêm yết, không phải giá giao dịch thực tế.

\section{Chỉ số và phương pháp phân tích}
\subsection{Bộ KPI sử dụng}
Các KPI được tính trên dữ liệu sau QA, kèm mẫu số và phạm vi. Với mẫu số bằng 0, kết quả được ghi là không xác định thay vì gán bằng 0.

\begin{table}[htbp]
\centering\small
\caption{Các chỉ số định lượng chính trong phạm vi nhóm mẫu.}\label{tab:kpi}
\begin{tabularx}{\linewidth}{@{}lYY@{}}
\toprule
KPI & Định nghĩa & Kết quả đã có căn cứ\\
\midrule
Nguồn cung quan sát & Số \code{listing\_id} duy nhất sau kiểm soát chất lượng ở mỗi mốc & 1.000 chỗ ở tại mỗi mốc; độ phủ nhóm mẫu đạt 100\%.\\
Mẫu giá hợp lệ & Số chỗ ở có giá dương và chuyển đổi được sang dạng số tại mốc tính giá & 980/1.000 chỗ ở tại 2026-06-15.\\
Giá trung vị & Trung vị giá niêm yết hợp lệ tại mốc tính giá & 213,55 CAD/đêm (n = 980).\\
Tỷ trọng toàn bộ nhà/căn hộ & Chỗ ở có loại \code{Entire home/apt} chia cho số chỗ ở có loại phòng hợp lệ & 87,5\% trong nhóm mẫu.\\
Tỷ trọng chỗ ở thuộc chủ nhà nhiều chỗ ở & Chỗ ở thuộc chủ nhà có hơn một chỗ ở trong nhóm mẫu, chia cho số chỗ ở có mã chủ nhà hợp lệ & 11,6\% tại mốc cuối; không đại diện cho các chỗ ở ngoài nhóm mẫu.\\
Thay đổi chủ nhà giữa các mốc & Số chỗ ở được giữ lại nhưng đổi \code{host\_id}, chia cho số cặp so sánh hợp lệ & 3/1.000 (0,3\%), 35/1.000 (3,5\%), 37/1.000 (3,7\%).\\
Độ phủ lịch 90 ngày & Số ngày--chỗ ở có trạng thái hợp lệ chia cho $1.000\times 90$ ngày & 99,78\%; 99,38\%; 100\%; 100\% theo bốn mốc dữ liệu.\\
\bottomrule
\end{tabularx}
\end{table}

Với vùng $g$ và mốc dữ liệu $t$, nguồn cung là $N_{g,t}=|L_{g,t}|$; giá trung vị được tính trên tập giá hợp lệ $P_{g,t}$ tại mốc cuối. Tỷ trọng toàn bộ nhà/căn hộ và tỷ trọng chủ nhà nhiều chỗ ở dùng mẫu số là các dòng có trường tương ứng hợp lệ. Tỷ lệ lịch còn sẵn (availability) được tính bằng số ngày có trạng thái mở chia cho số ngày có trạng thái quan sát được (mở hoặc không mở); độ phủ lịch là số ngày quan sát được chia cho số ngày kỳ vọng. \textbf{Tỷ lệ lịch còn sẵn không phải tỷ lệ lấp đầy thực tế.}

\subsection{Thiết kế phân tích bài đánh giá và mô hình ngôn ngữ lớn (LLM)}
Các bài đánh giá được loại trùng theo \code{review\_id} giữa các mốc dữ liệu. Bộ 100 nhãn tay được gán độc lập theo hướng dẫn chung, mỗi thành viên gán 20 bài đánh giá. Phân bố nhãn trong tệp CSV đã chốt là 90 positive (tích cực), 7 mixed (hỗn hợp) và 3 negative (tiêu cực). Phân bố nghiêng mạnh về nhãn positive (tích cực); đây là đặc điểm của tập nhãn được giao, không thể suy rộng thành tỷ lệ hài lòng của toàn bộ khách Airbnb Toronto.

Phương pháp đường cơ sở sử dụng từ khóa đa ngữ, ngữ cảnh phủ định, liên từ đối lập và một số mẫu diễn đạt sự cố nhỏ đã được xử lý. Cấu hình Gemini hiện sử dụng mô hình \code{gemini-3.5-flash-lite}, nhiệt độ 0, đầu ra JSON và tám khía cạnh: location, cleanliness, host, value, amenities, noise, safety và check-in. Bộ nhớ đệm lưu mã bài đánh giá, mã băm văn bản, tên mô hình, phiên bản câu lệnh và danh mục khía cạnh để ngăn dùng nhầm dự đoán khi đầu vào thay đổi.

\begin{equation*}
\text{Độ chính xác}=\frac{\text{số dự đoán đúng}}{n},\qquad
F_{1,c}=\frac{2P_cR_c}{P_c+R_c},\qquad
F_{1,\mathrm{macro}}=\frac{1}{K}\sum_{c=1}^{K}F_{1,c}.
\end{equation*}

Phương pháp đường cơ sở và Gemini cần được đánh giá trên cùng 100 \code{review\_id} có nhãn tay; báo cáo cần kèm ma trận nhầm lẫn, F1 theo từng lớp và ví dụ dự đoán sai. Tại thời điểm kiểm tra repo, bộ nhớ đệm Gemini thật cho 1.000 bài đánh giá chưa được tạo. Vì vậy, quy trình đầy đủ dừng khi thiếu bộ nhớ đệm và chưa có chỉ số độ chính xác/F1 vĩ mô của Gemini đã được xác minh. Không sử dụng bộ nhớ đệm giả trong kiểm thử hoặc số liệu của bản nháp cũ của TV5 làm kết quả thật.

\section{Kết quả và phân tích}
\subsection{RQ1: Cấu trúc mẫu được phân bổ trên tám vùng phân tích}
Tám vùng A01--A08 là các hình chữ nhật phân tích được xác định trong bảng tính phân công. Số chỗ ở được giao ở mỗi vùng dao động từ 120 đến 130; tổng cộng đúng 1.000 chỗ ở. Hình~\ref{fig:areas} mô tả quy mô mẫu chứ không phải mật độ toàn thị trường.

\begin{figure}[htbp]
\centering
\begin{tikzpicture}
\begin{axis}[
  ybar,bar width=13pt,width=0.88\linewidth,height=5.0cm,
  ymin=0,ymax=145,ylabel={Số chỗ ở trong nhóm mẫu},
  symbolic x coords={A01,A02,A03,A04,A05,A06,A07,A08},xtick=data,
  x tick label style={font=\small},ymajorgrids=true,grid style={gray!20},
  nodes near coords,nodes near coords style={font=\scriptsize},
  title={Quy mô nhóm mẫu theo vùng phân tích}]
\addplot[fill=blue!35,draw=blue!60!black] coordinates
 {(A01,130)(A02,130)(A03,126)(A04,121)(A05,120)(A06,121)(A07,130)(A08,122)};
\end{axis}
\end{tikzpicture}
\caption{Phân bổ 1.000 chỗ ở theo tám vùng A01--A08; số liệu lấy từ tệp kê khai phân công của nhóm.}
\label{fig:areas}
\end{figure}

\needspace{8\baselineskip}
Trong nhóm mẫu, 87,5\% chỗ ở có loại phòng \code{Entire home/apt}; tỷ lệ này không đại diện cho toàn Toronto. Ở mốc 2026-06-15, có 980 chỗ ở có giá hợp lệ và giá trung vị là 213,55 CAD/đêm. Dù quy trình xuất KPI theo tám vùng và 125 khu vực hành chính, bảng/biểu đồ giá của lần chạy đầy đủ gần nhất chưa được lưu trong kho mã; vì vậy, chưa đủ căn cứ xếp hạng khu vực đắt/rẻ. \pending{TV3/TV1 bổ sung bảng và biểu đồ không gian từ lần chạy chốt, kèm số quan sát và độ phủ.}

\subsection{RQ2: Số lượng chỗ ở trong nhóm mẫu ổn định nhưng mã chủ nhà có thay đổi}
Mỗi mốc dữ liệu đều có đủ 1.000 chỗ ở trong danh sách được giao; không ghi nhận mã mới hoặc vắng mặt trong ba khoảng so sánh liền kề. Tuy nhiên, mã chủ nhà của cùng một chỗ ở có thay đổi. Số cặp thay đổi lần lượt là 3, 35 và 37 trên 1.000 cặp có thể so sánh; loại phòng không thay đổi trong các cặp được báo cáo.

\begin{figure}[htbp]
\centering
\begin{tikzpicture}
\begin{axis}[
  ybar,bar width=22pt,width=0.72\linewidth,height=4.5cm,
  ymin=0,ymax=4.5,ylabel={Chỗ ở đổi mã chủ nhà (\%)},
  symbolic x coords={P1,P2,P3},xtick=data,
  xticklabels={09--12/2025,12/2025--03/2026,03--06/2026},
  x tick label style={font=\scriptsize},ymajorgrids=true,grid style={gray!20},
  nodes near coords={\pgfmathprintnumber[fixed,precision=1,use comma]{\pgfplotspointmeta}},nodes near coords style={font=\small},
  title={Tỷ lệ thay đổi mã chủ nhà giữa hai mốc liền kề}]
\addplot[fill=orange!45,draw=orange!70!black] coordinates
 {(P1,0.3)(P2,3.5)(P3,3.7)};
\end{axis}
\end{tikzpicture}
\caption{Tỷ lệ mã chủ nhà thay đổi trên 1.000 chỗ ở được giữ lại ở mỗi khoảng.}
\label{fig:host}
\end{figure}

Mức thay đổi tăng từ 0,3\% ở khoảng đầu lên 3,5\% và 3,7\% ở hai khoảng sau. Đây là dấu hiệu cần thận trọng khi nối KPI chủ nhà qua thời gian: tỷ lệ chủ nhà nhiều chỗ ở nên được tính theo từng mốc dữ liệu, không giả định mã chủ nhà bất biến. Việc mã chủ nhà đổi có thể do cập nhật tài khoản hoặc ghi nhận dữ liệu; dữ liệu hiện có không chứng minh chuyển quyền sở hữu hay nguyên nhân pháp lý.

Đối với bảng lịch, mẫu số kỳ vọng của cửa sổ 90 ngày là 90.000 ngày--chỗ ở tại mỗi mốc. Số ngày thiếu lần lượt là 200, 560, 0 và 0; tương ứng độ phủ 99,78\%, 99,38\%, 100\% và 100\%.

\begin{figure}[htbp]
\centering
\begin{tikzpicture}
\begin{axis}[
  ybar,bar width=22pt,width=0.72\linewidth,height=4.4cm,
  ymin=99,ymax=100.15,ylabel={Độ phủ quan sát (\%)},
  symbolic x coords={D1,D2,D3,D4},xtick=data,
  xticklabels={09/2025,12/2025,03/2026,06/2026},
  x tick label style={font=\scriptsize},ymajorgrids=true,grid style={gray!20},
  yticklabel={\pgfmathprintnumber[fixed,precision=1,use comma]{\tick}},
  nodes near coords={\pgfmathprintnumber[fixed,precision=2,use comma]{\pgfplotspointmeta}},nodes near coords style={font=\scriptsize},
  title={Độ phủ dữ liệu lịch trong cửa sổ 90 ngày}]
\addplot[fill=teal!35,draw=teal!65!black] coordinates
 {(D1,99.78)(D2,99.38)(D3,100)(D4,100)};
\end{axis}
\end{tikzpicture}
\caption{Độ phủ = số ngày--chỗ ở quan sát được / 90.000 ngày--chỗ ở kỳ vọng; không phải tỷ lệ lấp đầy.}
\label{fig:calendar}
\end{figure}

Độ phủ cao cho thấy dữ liệu lịch gần đầy đủ trong cửa sổ phân tích, nhưng trạng thái \code{available = f} gộp cả ngày đã được đặt và ngày chủ nhà chủ động khóa. Do đó, không thể dùng $1-\text{tỷ lệ sẵn có}$ để suy ra tỷ lệ lấp đầy thực tế. Các giá trị tỷ lệ lịch còn sẵn trong 30/90 ngày và phân tách ngày thường/cuối tuần cần được xuất lại từ đúng lần chạy quy trình cuối trước khi đưa vào bản nộp. \pending{TV4 xuất và đối chiếu tỷ lệ lịch còn sẵn 30/90 ngày; không thay chỉ số này bằng độ phủ quan sát.}

\subsection{RQ3: Bộ nhãn tay đã hoàn thành, nhưng đánh giá mô hình ngôn ngữ lớn còn chờ}
Bộ nhãn thủ công có đủ 100 \code{review\_id} được đánh dấu trong bảng tính phân công, gồm 90 nhãn positive (tích cực), 7 nhãn mixed (hỗn hợp) và 3 nhãn negative (tiêu cực). Phân bố này cho thấy tập kiểm chứng có mất cân bằng lớp; F1 vĩ mô (Macro-F1) và F1 theo từng lớp cần được báo cáo song song với độ chính xác để tránh chỉ số tổng thể che khuất hiệu năng trên hai lớp mixed và negative.

\begin{table}[htbp]
\centering\small
\caption{Phân bố 100 nhãn cảm xúc thủ công đã chốt.}\label{tab:labels}
\begin{tabular}{@{}lrr@{}}
\toprule
Nhãn & Số bài đánh giá & Tỷ lệ\\
\midrule
Positive & 90 & 90,0\%\\
Mixed & 7 & 7,0\%\\
Negative & 3 & 3,0\%\\
\midrule
Tổng & 100 & 100,0\%\\
\bottomrule
\end{tabular}
\end{table}

Kho mã có phương pháp đường cơ sở dựa trên từ khóa, mã tích hợp và kiểm tra Gemini, cơ chế bộ nhớ đệm có lưu thông tin nguồn gốc và hàm đánh giá dự đoán so với nhãn chuẩn. Tuy nhiên, ở phiên bản mã nguồn nêu đầu báo cáo, bộ nhớ đệm Gemini thật cho 1.000 bài đánh giá chưa có; chưa có bảng so sánh mô hình ngôn ngữ lớn với phương pháp đường cơ sở trên 100 nhãn đã chốt. Vì thế, các tỷ lệ khen/chê theo tám khía cạnh và chỉ số trải nghiệm khách lưu trú chưa được xem là kết quả thực nghiệm. \pending{TV5 tạo bộ nhớ đệm thật và kiểm tra đủ 1.000 mã bài đánh giá; chạy quy trình đầy đủ, bổ sung độ chính xác, F1 vĩ mô, phân tích lỗi, số lượng token và chi phí thực tế hoặc ước tính có căn cứ.}

\section{Thảo luận và kết luận}
Kết quả hiện có cho thấy ba điểm chính. Thứ nhất, quy trình xử lý giữ được toàn bộ 1.000 chỗ ở trong nhóm mẫu ở cả bốn mốc, và bước kiểm soát chất lượng không làm giảm số lượng chỗ ở; các trường hợp thiếu giá ở mốc cuối vẫn được giữ trong nhóm nhưng không tham gia phép tính giá. Thứ hai, mã chủ nhà thay đổi nhiều hơn ở hai khoảng sau, vì vậy các chỉ số chủ nhà cần được tính riêng theo từng mốc và không được diễn giải thành thay đổi chủ sở hữu. Thứ ba, độ phủ dữ liệu lịch 90 ngày ở mức 99,38--100\%, nhưng độ phủ lịch không đồng nghĩa với công suất khai thác thực tế.

Với bên nhận phân tích giả định, các chỉ số này phù hợp để theo dõi quy mô nhóm mẫu, cơ cấu loại phòng, giá niêm yết cuối kỳ và chất lượng quan sát. Trước khi đưa ra khuyến nghị về khác biệt giữa khu vực hoặc trải nghiệm khách, nhóm cần hoàn tất bảng KPI/hình theo vùng từ lần chạy cuối và đánh giá Gemini trên bộ nhãn tay. Không có căn cứ trong dữ liệu hiện tại để suy ra doanh thu thực, tỷ lệ lấp đầy, nhu cầu du lịch hoặc quan hệ nhân quả.

\section{Hạn chế}
\begin{enumerate}
  \item \textbf{Phạm vi nhóm mẫu:} Kết quả chỉ mô tả 1.000 mã chỗ ở và 5.000 mã bài đánh giá được giao; không đại diện toàn bộ Airbnb tại Toronto.
  \item \textbf{Các mốc dữ liệu:} Bốn mốc rời rạc không mô tả liên tục mọi biến động. Ngày 2025-12-12 trong bảng tính phân công khác 2025-11-11 trên trang đề; cần xác nhận trước khi nộp \cite{project}.
  \item \textbf{Giá:} Chỉ dùng giá trong bảng listings tại 2026-06-15; 20 chỗ ở thiếu giá và giá trung vị chỉ đại diện 980 chỗ ở có giá hợp lệ. Giá niêm yết không phải giá giao dịch.
  \item \textbf{Không gian:} A01--A08 là vùng chữ nhật theo bảng tính phân công, không phải đơn vị hành chính. Phân tích theo khu vực hành chính phải báo riêng độ phủ và số quan sát.
  \item \textbf{Dữ liệu lịch:} Ngày không còn sẵn có thể đã được đặt hoặc bị chủ nhà khóa; tỷ lệ lịch còn sẵn không thay thế tỷ lệ lấp đầy.
  \item \textbf{Bài đánh giá và mô hình ngôn ngữ lớn:} Bài đánh giá có thiên lệch tự chọn; 100 nhãn tay mất cân bằng lớp. Bộ nhớ đệm Gemini và đánh giá mô hình chưa hoàn tất, nên chưa thể kết luận về chất lượng phân loại cảm xúc hoặc trích xuất khía cạnh.
\end{enumerate}

\section{Khả năng tái lập, minh bạch AI và đóng góp}
\textbf{Môi trường.} Kho mã ghi nhận Python 3.13.7, các thư viện được cố định phiên bản trong \code{requirements.txt}; bộ kiểm thử hiện tại có 139 phép thử đạt. Lệnh chạy phần dữ liệu thật không dùng mô hình ngôn ngữ lớn:
\begin{quote}\small\ttfamily
python run\_pipeline.py --config configs/toronto.yaml --skip-llm --no-download
\end{quote}
Lệnh chạy toàn bộ quy trình sau khi có bộ nhớ đệm:
\begin{quote}\small\ttfamily
python run\_pipeline.py --config configs/toronto.yaml --no-download
\end{quote}
Lệnh thứ hai hiện chưa thể hoàn tất do chưa có bộ nhớ đệm Gemini. Kiểm thử bằng dữ liệu mẫu và bộ nhớ đệm giả chỉ xác nhận hoạt động của phần mềm, không chứng minh chất lượng dự đoán của Gemini. Nhóm chưa tạo thẻ phiên bản \code{final}; mã SHA nêu ở đầu báo cáo chỉ xác định phiên bản mã nguồn dùng để đối chiếu bản thảo.

\textbf{Minh bạch AI.} Theo \code{AI\_USAGE.md}, OpenAI Codex được dùng để hỗ trợ dựng khung kho mã, tích hợp, kiểm tra và biên tập tài liệu/mã. Gemini được thiết kế cho bước phân tích bài đánh giá nhưng chưa được gọi trong lần kiểm chứng dữ liệu thật gần nhất và chưa có bộ nhớ đệm kết quả thật. 100 nhãn cảm xúc thủ công được con người gán và giữ nguyên; không dùng AI để tạo nhãn chuẩn. Các số liệu trong báo cáo được đối chiếu với tệp README, cấu hình, tệp kê khai, báo cáo kiểm chứng và phần bàn giao của các thành viên \cite{repo}.

\textbf{Đóng góp.} TV1 -- Dương Nhật Huy: điều phối quy trình xử lý, phân tích nhóm mẫu và biến động thuộc tính; TV2 -- Mai Anh Đức: kiểm soát chất lượng, phân tích khám phá và đánh giá tác động chất lượng dữ liệu; TV3 -- Nguyễn Quốc Hiếu: chỉ số nguồn cung/giá/loại phòng và phân tích không gian; TV4 -- Vũ Phạm Khánh Nam: xử lý dữ liệu lịch và tỷ lệ sẵn có; TV5 -- Phan Ngọc Huy: xử lý bài đánh giá, phương pháp đường cơ sở, mô hình ngôn ngữ lớn, đánh giá và chỉ số trải nghiệm khách. \pending{các thành viên xác nhận lại mô tả đóng góp và bổ sung MSSV trước khi nộp.}

\begin{thebibliography}{9}\footnotesize\sloppy
\setlength{\itemsep}{0pt}
\bibitem{inside} Inside Airbnb. \textit{Get the Data}. Dữ liệu Toronto, Canada, giấy phép CC BY 4.0. \href{https://insideairbnb.com/get-the-data/}{Trang dữ liệu}. Truy cập ngày 17/11/2026.
\bibitem{snap1} Inside Airbnb. Dữ liệu Toronto, mốc 2025-09-11. \href{https://data.insideairbnb.com/canada/on/toronto/2025-09-11/}{Trang dữ liệu mốc 2025-09-11}.
\bibitem{snap2} Inside Airbnb. Dữ liệu Toronto, mốc 2025-12-12 theo bảng tính phân công nhóm (chờ xác nhận). \href{https://data.insideairbnb.com/canada/on/toronto/2025-12-12/}{Trang dữ liệu mốc 2025-12-12}.
\bibitem{snap3} Inside Airbnb. Dữ liệu Toronto, mốc 2026-03-19. \href{https://data.insideairbnb.com/canada/on/toronto/2026-03-19/}{Trang dữ liệu mốc 2026-03-19}.
\bibitem{snap4} Inside Airbnb. Dữ liệu Toronto, mốc 2026-06-15. \href{https://data.insideairbnb.com/canada/on/toronto/2026-06-15/}{Trang dữ liệu mốc 2026-06-15}.
\bibitem{project} Viện Trí tuệ nhân tạo, Trường Đại học Công nghệ, ĐHQGHN. \textit{Bài tập lớn Inside Airbnb -- UET.DSE2049, năm học 2026--2027}. \href{https://courses.iaidev.com/programming-for-data-processing/2627-1/projects/project_airbnb.html}{Trang đề bài}.
\bibitem{repo} Nhóm 10. \textit{Mã nguồn, cấu hình và hồ sơ kiểm chứng dự án Inside Airbnb Toronto}. GitHub, nhánh \code{main}, commit \code{d045d55869b76c21b92b8b48bf2f0417a929efca}. \href{https://github.com/uet-iai-notebook-labs-2026/2627-LTXLDL-Airbnb-dse2049_1-nhom10}{Kho mã Nhóm 10}.
\end{thebibliography}

\vspace{2pt}
% Checklist biên tập nội bộ: hoàn tất bộ nhớ đệm Gemini và kết quả đánh giá; cập nhật các bảng/hình; xác nhận mốc 2025-12-12; bổ sung MSSV; cập nhật mã SHA.

\end{document}
